La géométrie du collège se démontre ici dans un espace euclidien quelconque, sans repère :
les points sont des vecteurs, la distance est la norme de leur différence, l'angle est celui
que forment deux vecteurs. Ce choix a un prix — aucune figure ne vient soutenir le
raisonnement — et un avantage : chaque énoncé vaut tel quel dans le plan comme dans
l'espace, et rien ne peut s'appuyer sur un dessin bien orienté.

Le théorème de Pythagore et sa réciproque se ramènent alors à une seule identité, celle qui
développe le carré de la norme d'une somme ; le théorème de Thalès, à la conservation des
rapports par homothétie. Les caractérisations — médiatrice, bissectrice, parallélogramme —
sont des équivalences, et chacune se démontre dans les deux sens, ce que la rédaction
scolaire escamote souvent.

Deux points demandent une attention particulière. La \emph{réciproque du parallélogramme}
est fausse sans l'hypothèse que le quadrilatère n'est pas croisé : le chapitre l'énonce avec
cette hypothèse et donne le contre-exemple qui montre qu'elle est nécessaire. Et les angles
\emph{orientés} vivent dans un fichier séparé, parce qu'ils exigent que le plan soit de
dimension deux — une hypothèse qui, laissée dans le fichier principal, contaminerait tous
les autres énoncés.
